\documentclass[10pt,a4paper]{amsart}
\usepackage[margin=19mm]{geometry}
\usepackage{amsmath,amssymb,mathtools}
\usepackage[scr=rsfs,cal=euler]{mathalfa}
\usepackage{xfrac}
\usepackage[unicode,hidelinks,bookmarks=false]{hyperref}
\newcommand{\ie}{i.e.\ }
\newcommand{\cf}{cf.\ }
\newcommand{\resp}{respectively\ }
\newcommand{\Subsection}{\subsection}
\providecommand{\nobreakdash}{\nobreak\hbox{-}}
\setlength{\emergencystretch}{2em}
\multlinegap0pt
\usepackage{booktabs}
\newtheorem{theorem}{Theorem}
\newtheorem{proposition}{Proposition}
\newtheorem{lemma}{Lemma}
\newtheorem{corollary}{Corollary}
\newtheorem{definition}{Definition}
\newtheorem{remark}{Remark}
\begin{document}
\title[Asymptotic Signal Subspace Recovery in Softmax Attention Mod]{Asymptotic Signal Subspace Recovery in Softmax Attention Models}
\author{Lan V. Truong; Faculty of Computer Science and Engineering}
\date{}
\maketitle
\noindent\emph{Source and licence.} Asymptotic Signal Subspace Recovery in Softmax Attention Models. Version arXiv:2606.22406v2 (CC BY 4.0). This material is distributed under the Creative Commons Attribution 4.0 International licence, http://creativecommons.org/licenses/by/4.0/, as recorded in the arXiv OAI licence metadata for this exact version and shown on the abstract page https://arxiv.org/abs/2606.22406v2. Full rights notice: licenses/arxiv-2606.22406-CC-BY-4.0.txt.\par\smallskip
\noindent\textit{Editorial scope.} Counted unit: the original \texttt{theorem} environment \texttt{thm:main} at source lines 376--385 together with its complete proof at lines 387--498; the delivered document keeps source lines 218--499 so that every referenced label and every citation resolves inside it. Native editable source is the paper's own concatenated TeX; the AMS wrapper is editorial formatting only. No publisher class, upstream script or unrelated artwork is redistributed.
\begin{lemma}[Smoothness of the Population Objective] \label{lem:lem1}
The unsupervised population function satisfies $J_{d}(q) \in C^{\infty}(S^{d-1})$.
\end{lemma}

\noindent Due to the isotropic profile of the random noise vectors, the high-dimensional objective collapses into a compact system determined solely by the latent axis projection.

\begin{lemma}[Rotational Invariance and Dimension Reduction] \label{lem:lem2}
There exists a smooth function $\Psi_{d}:[-1,1] \rightarrow \mathbb{R}$ such that $J_{d}(q) = \Psi_{d}(\langle q, \xi_{d} \rangle)$.
\end{lemma}

\noindent To measure the concentration of individual token weights under this representation, we provide a uniform matrix norm envelope bound on the softmax allocations.

\begin{lemma}[Concentration of Attention Weights] \label{lem:lem3}
The squared Euclidean norm of the attention weights satisfies 
\begin{equation}
    \|a(q)\|^2 =O\left( \frac{L e^{\beta^2} }{R(L-R)}\right) \qquad\text{a.s.}
\end{equation}
\end{lemma}

\noindent We complete the landscape characterization by establishing that the reduced representation is strictly monotonic along the alignment vector field, preventing spurious local traps.

\begin{lemma}[Monotonicity of the Reduced Objective] \label{lem:lem4}
As $R \to \infty, L-R \to \infty, \beta=O(1), \theta_d =O(1)$, the objective satisfies $\Psi_{d}^{\prime}(\rho) > 0$ for all intermediate alignment values $\rho \in [-1,1]$.
\end{lemma}


\subsection{Algorithmic Gradients \& Tracking Dynamics}
With the underlying landscape characteristics defined, we next isolate individual stochastic updates and use multi-index Taylor expansions to transition the discrete steps into continuous paths.

\begin{lemma}[Uniform Gradient Variance Bound] \label{lem:lem5}
Under the condition that the inverse temperature and signal strength parameters satisfy $\beta =O(1), \theta_d = O(1)$, the unsupervised stochastic gradient satisfies a dimension-dependent second moment condition:
\begin{equation}
    \mathbb{E}\left[ \|g_k\|^{2} \right] \le C d^{4}.
\end{equation}
\end{lemma}

\noindent By projecting the empirical gradient variance onto the tangent plane of the hypersphere, we decompose the discrete recurrence update into a true drift path and an optimization error envelope.

\begin{lemma}[Discrete Taylor Trajectory Projections] \label{lem:lem6}
The projected SGD updating equations expand as:
\begin{equation}
    q_{k+1} = q_k + \eta_k(I_d - q_k q_k^{\top})g_k + r_k,
\end{equation}
where the cumulative remainder vector satisfies the uniform bound $\|r_k\| \le C \eta_k^{2} \|g_k\|^{2}$.
\end{lemma}


\subsection{Stochastic Convergence \& Martingale Bounds}
The final phase of our auxiliary analysis provides the requisite martingale limits and Lipschitz tracking metrics required to execute the continuous ODE method.

\begin{lemma}[Global Drift Field Lipschitz Continuity] \label{lem:lem7}
The continuous drift vector field $H_{d}(q) = (I_d - q q^\top)\nabla_q J_d(q)$ is globally Lipschitz continuous on the hypersphere $S^{d-1}$.
\end{lemma}

\noindent To guarantee that empirical noise fluctuations balance out uniformly as iterations accumulate, we establish the convergence properties of the noise filtration array.

\begin{lemma}[Martingale Array Convergence Bounds] \label{lem:lem8}
The stochastically accumulating martingale sum series $\sum_{k=0}^{\infty} \eta_k  M_{k+1}$ converges almost surely.
\end{lemma}

\noindent Finally, we show that the accumulated mathematical approximations introduced during the Taylor projections vanish over time.

\begin{lemma}[Remainder Convergent Summability] \label{lem:lem9}
The cumulative Taylor remainder sequence is absolutely summable almost surely: 
\begin{equation}
    \sum_{k=0}^{\infty} \|r_k\| < \infty.
\end{equation}
\end{lemma}
\subsection{Time-Scale Mapping and Continuous Interpolation}
To study the discrete-time sequence $\{q_k\}_{k=0}^{\infty}$ using dynamical systems tools, we map the discrete iteration counter $k$ onto a continuous timeline. We define the cumulative time intervals $\tau_0 = 0$ and $\tau_k = \sum_{l=0}^{k-1} \eta_l$ for $k \ge 1$, where $\eta_l$ is the step size. 

Let $m(t) = \max\{k \ge 0 : \tau_k \le t\}$ denote the discrete step counter corresponding to the continuous time parameter $t$. We then construct the continuous-time piecewise constant interpolation trajectory $\bar{q}(t)$ defined by:
\begin{equation}
    \bar{q}(t) = q_{m(t)}, \quad t \in [0, \infty).
\end{equation}
Throughout the subsequent analysis, discrete-time iteration updates are denoted by a subscript index ($q_k$), whereas continuous-time trajectory coordinates and solutions to the limit differential equation $\dot{q} = H_d(q)$ are designated as functions of time ($q(t)$).
\subsection{Asymptotic Trajectory Tracking Proofs}

\begin{proposition}[Asymptotic Pseudo-Trajectory Approximation] \label{prop:tracking}
Let $q(t)$ be the continuous interpolation step trajectory. Under Lemmas~\ref{lem:lem1} through \ref{lem:lem9}, the discrete updates tracking sequence maps uniformly onto the continuous ODE system flow framework.
\end{proposition}

\begin{proof}[Proof of Proposition \ref{prop:tracking}]
Recall the one-step recursion established in Lemma~\ref{lem:lem6}:
\begin{align}
q_{k+1}= q_k + \eta_k H_d(q_k) + \eta_k M_{k+1} + r_k \label{eq16}.
\end{align}
To demonstrate convergence under the continuous framework, we partition our verification across the following structural stages:

\begin{description}
    \item[Step 1: Integral representation.] Fix $t \ge 0$ and $s \ge 0$. Summing the discrete recursion matching indices from $k = m(t)$ to $m(t+s)-1$ in \eqref{eq16} yields:
    \begin{align}
  \bar{q}(t+s)-\bar{q}(t)&=  q_{m(t+s)} - q_{m(t)} \nonumber \\
  &= \sum_{k=m(t)}^{m(t+s)-1} \eta_k H_d(q_k) + \sum_{k=m(t)}^{m(t+s)-1} \eta_k M_{k+1} + \sum_{k=m(t)}^{m(t+s)-1} r_k \label{eq17}.
    \end{align}
    
    \item[Step 2: Control of the interpolation error.] Let $k = m(u)$ for a given continuous time coordinate index. Evaluating the step offset bounds $\tau_k \le u < \tau_{k+1}$, the piecewise trajectory error satisfies:
    \begin{align}
    \| \bar{q}(u) - q_k \| = \| q_k - q_k \| = 0.
    \end{align}
    Consequently, we can write the main continuous tracking drift component directly as an integrated representation over the piecewise field:
    \begin{align}
    \sum_{k=m(t)}^{m(t+s)-1} \eta_k H_d(q_k) = \int_{t}^{t+s} H_d(\bar{q}(u)) \, du.
    \end{align}

    \item[Step 3: Evaluation of the continuous drift discrepancy.] To handle the difference between the actual trajectory path $q(t+u)$ and its step-wise representation $\bar{q}(t+u)$, we subtract the continuous differential equation component:
    \begin{align}
    q(t+s) - q(t) = \int_{t}^{t+s} H_d(q(u)) \, du \label{eq20}.
    \end{align}
    Taking the difference between the discrete interpolation series in~\eqref{eq17} and the continuous flow in~\eqref{eq20}  yields:
    \begin{align}
    \bar{q}(t+s) - q(t+s) = \bar{q}(t) - q(t) + \int_{t}^{t+s} \left[ H_d(\bar{q}(u)) - H_d(q(u)) \right] \, du + E(t,s)
    \end{align}
    where the combined error envelope tracking array maps directly to:
    \begin{align}
    E(t,s) = \sum_{k=m(t)}^{m(t+s)-1} \eta_k M_{k+1} + \sum_{k=m(t)}^{m(t+s)-1} r_k.
    \end{align}

    \item[Step 4: Application of Martingale and Remainder limits.] By Lemma~\ref{lem:lem8}, the martingale tracking array sequence $\sum_{k=0}^{\infty} \eta_k M_{k+1}$ converges almost surely. By Cauchy uniformity \cite{Royden}, this implies:
    \begin{align}
    \lim_{t \to \infty} \sup_{0 \le s \le T} \left\| \sum_{k=m(t)}^{m(t+s)-1} \eta_k M_{k+1} \right\| = 0 \qquad \text{a.s.}
    \end{align}
    Similarly, applying the absolute summability property from Lemma~\ref{lem:lem9}, the trailing remainder norm satisfies:
    \begin{align}
    \lim_{t \to \infty} \sup_{0 \le s \le T} \left\| \sum_{k=m(t)}^{m(t+s)-1} r_k \right\| \le \lim_{t \to \infty} \sum_{k=m(t)}^{\infty} \|r_k\| = 0 \qquad \text{a.s.}
    \end{align}
    Combining these limits under the supremum metric over a compact tracking window gives $\lim_{t \to \infty} \sup_{0 \le s \le T} \|E(t,s)\| = 0$ almost surely.

    \item[Step 5: Invoking the Lipschitz bound via Grönwall's Inequality.] Let $e(t) = \sup_{0 \le s \le T} \|\bar{q}(t+s) - q(t+s)\|$. Using the global Lipschitz constant $L$ proven in Lemma~\ref{lem:lem7}:
    \begin{align}
    \|\bar{q}(t+s) - q(t+s)\| \le \|\bar{q}(t) - q(t)\| + L \int_{t}^{t+s} \|\bar{q}(u) - q(u)\| \, du + \|E(t,s)\|.
    \end{align}
    Applying Grönwall's inequality~\cite{hirsch2013differential} yields:
    \begin{align}
    \sup_{0 \le s \le T} \|\bar{q}(t+s) - q(t+s)\| \le \left( \|\bar{q}(t) - q(t)\| + \sup_{0 \le s \le T} \|E(t,s)\| \right) e^{LT}.
    \end{align}

    \item[Step 6: Final asymptotic trajectory matching proof.] Taking the limit inferior and superior over time, the initialization offset $\|\bar{q}(t) - q(t)\|$ goes to 0 by definition of the matched start coordinates. Because $\lim_{t \to \infty} \sup_{0 \le s \le T} \|E(t,s)\| = 0$, we have:
    \begin{align}
    \lim_{t \to \infty} \sup_{0 \le s \le T} \|\bar{q}(t+s) - q(t+s)\| = 0 \qquad \text{a.s.}
    \end{align}
\end{description}
This completes the verification that the discrete-time optimization processes match the continuous-time ordinary differential equation limits almost surely.
\end{proof}

\section{Main Convergence Results}

With the static geometry of the landscape and the asymptotic tracking capabilities established in Section 5, we are now positioned to state our main convergence guarantees. We begin by formalizing the definition of the limit set of the continuous-time dynamical system, which acts as the target for our discrete stochastic trajectory.

\begin{definition}[Asymptotic Equilibrium Set] \label{def:equilibrium_set}
The equilibrium target set $\Lambda \subset S^{d-1}$ for the continuous projected field flow $H_d(q) = (I_d - qq^\top)\nabla_q J_d(q)$ is defined as the set of vanishing drift configurations:
\begin{equation}
    \Lambda = \left\{ q \in S^{d-1} : H_d(q) = 0 \right\}.
\end{equation}
\end{definition}

\noindent By leveraging the monotonicity property established in Lemma~\ref{lem:lem4}, we prove that this equilibrium set consists exclusively of the true latent signal direction and its exact inverse, showing that no spurious local maxima or saddle points can attract the optimization path inside the sphere.


\begin{theorem}[Global Convergence to the Signal Subspace] \label{thm:main}
Let $\{q_k\}_{k=0}^{\infty}$ be the sequence of query vectors generated by the projected stochastic gradient ascent update scheme on the hypersphere $S^{d-1}$. Assume the step-size sequence satisfies the standard Robbins-Monro conditions:
\begin{equation}
    \sum_{k=0}^{\infty} \eta_k = \infty \quad \text{and} \quad \sum_{k=0}^{\infty} \eta_k^2 < \infty.
\end{equation}
Furthermore, assume the system operates in the high-dimensional token scaling regime where $R \to \infty$, $L-R \to \infty$, and the inverse temperature and signal strength parameters satisfy $\beta =O(1), \theta_d = O(1)$. Then, the sequence of learned queries converges almost surely to the isolated latent signal subspace:
\begin{equation}
    \lim_{k \to \infty} \inf_{q^* \in \{\pm \xi_d\}} \|q_k - q^*\| = 0, \qquad \text{a.s.}
\end{equation}
\end{theorem}

\begin{proof}[Proof of Theorem~\ref{thm:main}]
Let
\begin{align*}
H_d(q)
=
(I_d - qq^\top)\nabla J_d(q)
\end{align*}
denote the limiting drift field, and consider the ODE
\begin{align}
\dot q = H_d(q).
\end{align}

By Lemma~\ref{lem:lem1}, $H_d$ is continuous on the compact manifold $S^{d-1}$.

We first characterize the equilibrium set of the ODE. Along any solution $q(t)$,
\begin{align*}
\frac{d}{dt}J_d(q(t))
&=
\left\langle \nabla J_d(q(t)), \dot q(t) \right\rangle \\
&=
\left\langle \nabla J_d(q(t)), (I_d - q(t)q(t)^\top)\nabla J_d(q(t)) \right\rangle \\
&=
\|(I_d - q(t)q(t)^\top)\nabla J_d(q(t))\|^2 \\
&=
\|H_d(q(t))\|^2 \ge 0.
\end{align*}

Moreover,
\begin{align}
\frac{d}{dt}J_d(q(t)) = 0
\quad \text{if and only if} \quad
H_d(q(t)) = 0.
\end{align}
Hence, $J_d$ is a strict Lyapunov function for the flow.

Now, by Lemma~\ref{lem:lem2} and Lemma~\ref{lem:lem4}, we have
\begin{align}
\nabla J_d(q)
=
\Psi_d'(\rho)\,\xi_d,
\qquad
\rho=\langle q,\xi_d\rangle,
\end{align}
and
\begin{align}
\Psi_d'(\rho) > 0,
\qquad
\rho \in [-1,1].
\end{align}

Hence, the equilibrium condition becomes
\begin{align}
\Psi_d'(\rho)\,(I_d - qq^\top)\xi_d = 0,
\end{align}
which implies
\begin{align}
(I_d - qq^\top)\xi_d = 0.
\end{align}

Therefore,
\begin{align}
\xi_d
=
qq^\top\xi_d
=
\langle q,\xi_d\rangle q
=
\rho q.
\end{align}

Taking norms gives
\begin{align}
1 = \|\xi_d\|
=
|\rho|\,\|q\|
=
|\rho|,
\end{align}
and therefore
\begin{align}
\rho = \pm 1.
\end{align}

Consequently, the equilibrium set is
\begin{align}
\Lambda = \{\xi_d, -\xi_d\}.
\end{align}

Since $J_d$ is a strict Lyapunov function, every internally chain transitive invariant set of the flow is contained in $\Lambda$~\cite{benaim1999dynamics}. 

Now, by Proposition~\ref{prop:tracking}, the affine interpolation of the stochastic iterates forms an asymptotic pseudo-trajectory of the ODE. Hence, by the Kushner–Clark theorem \cite{kushner1978, benaim1999dynamics}, the limit set $\omega(\{q_k\})$ is almost surely an internally chain transitive invariant set of the flow generated by $H_d$.

Since the only internally chain transitive invariant sets are the equilibria $\pm \xi_d$, it follows that
\begin{align}
\omega(\{q_k\})
\subseteq
\{\pm\xi_d\}
\qquad \text{a.s.}
\end{align}

Therefore,
\begin{align}
\lim_{k\to\infty}
\inf_{q^\ast\in\{\pm\xi_d\}}
\|q_k-q^\ast\|
=
0
\qquad \text{a.s.}
\end{align}

This completes the proof.
\end{proof}


\begin{thebibliography}{99}
\bibitem[Royden]{Royden} H.~Royden and P.~Fitzpatrick. \newblock {\em Real Analysis}. \newblock Pearson, 4th edition, 2010.
\bibitem[benaim1999dynamics]{benaim1999dynamics} Michel Bena{\"i}m. \newblock Dynamics of stochastic approximation algorithms. \newblock {\em S{\'e}minaire de Probabilit{\'e}s XXXIII}, 1709:1--68, 1999.
\bibitem[hirsch2013differential]{hirsch2013differential} Morris~W. Hirsch, Stephen Smale, and Robert~L. Devaney. \newblock {\em Differential Equations, Dynamical Systems, and an Introduction to Chaos}. \newblock Academic Press, 3rd edition, 2013.
\bibitem[kushner1978]{kushner1978} Harold~J. Kushner and Dean~S. Clark. \newblock {\em Stochastic Approximation Methods for Constrained and Unconstrained Systems}. \newblock Springer, 1978.
\end{thebibliography}
\end{document}
